\documentclass[12pt,a4paper]{article}
\usepackage[utf8]{vietnam}
\usepackage{amsmath,amssymb}
\usepackage{tabularx}
\usepackage{tikz}
\usepackage{xcolor}
\usepackage[top=0.6cm, bottom=0.6cm, left=0.6cm, right=0.6cm]{geometry}
\usepackage[dethi]{ex_test}

\setlength{\parskip}{0pt}
\setlength{\parindent}{0pt}

% Dòng kẻ chấm mờ đều đặn chuẩn viết tay (khoảng cách 20.5pt = ~7.2mm rộng rãi)
\newcommand{\dongke}[1]{%
  \begingroup%
  \setlength{\parskip}{0pt}%
  \baselineskip=20.5pt%
  \vspace{1.5mm}%
  \foreach \i in {1,...,#1}{%
    \noindent{\color{gray!60}\dotfill}\par%
  }%
  \vspace{1mm}%
  \endgroup%
}

% Định nghĩa ô bài tập tự luận trong bảng
\newcommand{\oBaiTap}[3][6]{%
  \vspace{1.5mm}%
  \noindent\textbf{#2} #3\par%
  \dongke{#1}%
}

% Câu trắc nghiệm 1 dòng (4 đáp án, mỗi đáp án 0.25\linewidth)
\newcommand{\cauTN}[6]{% \cauTN{số câu}{đề}{A}{B}{C}{D}
  \noindent\textbf{Câu #1.} #2\par\vspace{0.4mm}%
  \noindent%
  \makebox[0.25\linewidth][l]{\textbf{A.} #3}%
  \makebox[0.25\linewidth][l]{\textbf{B.} #4}%
  \makebox[0.25\linewidth][l]{\textbf{C.} #5}%
  \makebox[0.25\linewidth][l]{\textbf{D.} #6}\par\vspace{0.8mm}%
}

% Câu trắc nghiệm 2 dòng (mỗi dòng 2 đáp án, mỗi đáp án 0.5\linewidth)
\newcommand{\cauTNB}[6]{% \cauTNB{số câu}{đề}{A}{B}{C}{D}
  \noindent\textbf{Câu #1.} #2\par\vspace{0.4mm}%
  \noindent%
  \makebox[0.5\linewidth][l]{\textbf{A.} #3}%
  \makebox[0.5\linewidth][l]{\textbf{B.} #4}\par\vspace{0.4mm}%
  \noindent%
  \makebox[0.5\linewidth][l]{\textbf{C.} #5}%
  \makebox[0.5\linewidth][l]{\textbf{D.} #6}\par\vspace{0.8mm}%
}

\begin{document}
\pagestyle{empty}

% =============================================================================
% MỖI MÃ ĐỀ GỒM ĐÚNG 2 TRANG: MẶT TRƯỚC (TRANG 1/2) VÀ MẶT SAU (TRANG 2/2)
% =============================================================================

% ------------------------- MẶT TRƯỚC (TRANG 1/2) -----------------------------
\noindent
\begin{tabularx}{\linewidth}{|X|c|}
\hline
\textbf{PHIẾU KIỂM TRA ĐỊNH KỲ TOÁN <K> - TUẦN <X> - BUỔI <Y>} & \textbf{Mã đề: 101 - Trang 1/2} \\
\textit{Chủ đề: <TÊN BÀI HỌC>} & \textit{Thời gian: 45 phút} \\
\hline
\end{tabularx}

\vspace{1mm}

\noindent
\begin{tabularx}{\linewidth}{|X|}
\hline
\rule[-3.5mm]{0pt}{10mm}\textbf{Họ và tên học sinh:} \hfill \textbf{Mã HS:} \makebox[3.5cm]{\dotfill} \\
\hline
\end{tabularx}

\vspace{1mm}

% Ô ĐIỂM SỐ & LỜI NHẬN XÉT: BẮT BUỘC \rule{0pt}{1.8cm}
\noindent
\begin{tabularx}{\linewidth}{|p{3.2cm}|X|}
\hline
\centering\textbf{ĐIỂM SỐ} & \centering\textbf{LỜI NHẬN XÉT CỦA THẦY CÔ} \tabularnewline
\hline
\rule{0pt}{1.8cm} & \tabularnewline
\hline
\end{tabularx}

\vspace{2mm}

% TUYỆT ĐỐI KHÔNG ĐẶT BẢNG ĐIỀN ĐÁP ÁN TRẮC NGHIỆM
\noindent\textbf{\large PHẦN I. TRẮC NGHIỆM KHÁCH QUAN (2,5 điểm)} \hfill \textit{(Khoanh tròn đáp án đúng)}\par
\vspace{1mm}

% 5 câu trắc nghiệm
\cauTN{1}{Đề bài...}{A}{B}{C}{D}
\cauTN{2}{Đề bài...}{A}{B}{C}{D}
\cauTN{3}{Đề bài...}{A}{B}{C}{D}
\cauTN{4}{Đề bài...}{A}{B}{C}{D}
\cauTN{5}{Đề bài...}{A}{B}{C}{D}

\vspace{1.5mm}

\noindent\textbf{\large PHẦN II. TỰ LUẬN (7,5 điểm)}

\noindent\textbf{Câu 1 (2,5 điểm). Giải các hệ phương trình sau:}
\vspace{0.8mm}

% Tự luận Câu 1: Dùng \oBaiTap[16] cho bảng 1x2 để lấp đầy sát đáy trang 1
\noindent
\begin{tabularx}{\linewidth}{|X|X|}
\hline
\oBaiTap[16]{a)}{$\begin{cases} ax + by = c \\ dx + ey = f \end{cases}$} & 
\oBaiTap[16]{b)}{$\begin{cases} mx + ny = p \\ qx + ry = s \end{cases}$} \\
\hline
\end{tabularx}

\newpage

% ------------------------- MẶT SAU (TRANG 2/2) ------------------------------
\noindent
\begin{tabularx}{\linewidth}{|X|c|}
\hline
\textbf{PHẦN II. TỰ LUẬN (Tiếp theo) - TOÁN <K> TUẦN <X> BUỔI <Y>} & \textbf{Mã đề: 101 - Trang 2/2} \\
\hline
\end{tabularx}

\vspace{2.5mm}

\noindent\textbf{\large Câu 2 (2,5 điểm). Rút gọn các biểu thức phân thức chứa căn sau:}
\vspace{0.8mm}

% Tự luận Câu 2: Tối thiểu \oBaiTap[14]
\noindent
\begin{tabularx}{\linewidth}{|X|X|}
\hline
\oBaiTap[14]{a)}{$A = \dots$\par\vspace{0.5mm}\textit{(với $x \ge 0, x \ne \dots$)}} & 
\oBaiTap[14]{b)}{$B = \dots$\par\vspace{0.5mm}\textit{(với $x \ge 0, x \ne \dots$)}} \\
\hline
\end{tabularx}

\vspace{2.5mm}

\noindent\textbf{\large Câu 3 (2,5 điểm). Bài toán thực tế:}
\vspace{0.8mm}

% Tự luận Câu 3: Tối thiểu \dongke{15} để lấp đầy sát đáy trang 2
\noindent
\begin{tabularx}{\linewidth}{|X|}
\hline
\vspace{0.5mm}
Đề bài bài toán thực tế...\par
\vspace{0.5mm}
$a)$ Câu hỏi ý a...\par
$b)$ Câu hỏi ý b...\par
\vspace{0.4mm}
\noindent\textit{Bài giải:}\dongke{15}
\\ \hline
\end{tabularx}

\newpage
% LẶP LẠI CHO CÁC MÃ ĐỀ TIẾP THEO (102, 103, 104)

\end{document}
